\documentclass[10pt,a4paper]{amsart}
\usepackage{amsmath,amssymb,amsthm,graphicx}
% Retained preamble for the Stacks Project chapter morphisms.tex (Morphisms of Schemes).
% The chapter file pinned-original-source.tex carries only \input{preamble} on its first line, so no
% preamble text lives inside the chapter itself.  The Stacks Project uses one single shared preamble for
% every chapter; the snapshot stages it verbatim in sources/stacks-chow/preamble.tex.  The block below is
% copied mechanically, line for line, from that file's lines 14-18 (the Xy-pic setup) and lines 154-217
% (the native theorem declarations together with the mathematical macros).  Nothing was added, removed or
% redefined.  The wrapper class and the amsmath/amssymb/amsthm/graphicx base packages are supplied by the
% assembler, so the class/\externaldocument part of the shared preamble is not repeated here.
% For commutative diagrams we use Xy-pic
\usepackage[all]{xy}

% We use 2cell for 2-commutative diagrams.
\xyoption{2cell}
% Theorem environments.
%
\theoremstyle{plain}
\newtheorem{theorem}[subsection]{Theorem}
\newtheorem{proposition}[subsection]{Proposition}
\newtheorem{lemma}[subsection]{Lemma}

\theoremstyle{definition}
\newtheorem{definition}[subsection]{Definition}
\newtheorem{example}[subsection]{Example}
\newtheorem{exercise}[subsection]{Exercise}
\newtheorem{situation}[subsection]{Situation}

\theoremstyle{remark}
\newtheorem{remark}[subsection]{Remark}
\newtheorem{remarks}[subsection]{Remarks}

\numberwithin{equation}{subsection}

% Macros
%
\def\lim{\mathop{\mathrm{lim}}\nolimits}
\def\colim{\mathop{\mathrm{colim}}\nolimits}
\def\Spec{\mathop{\mathrm{Spec}}}
\def\Hom{\mathop{\mathrm{Hom}}\nolimits}
\def\Ext{\mathop{\mathrm{Ext}}\nolimits}
\def\SheafHom{\mathop{\mathcal{H}\!\mathit{om}}\nolimits}
\def\SheafExt{\mathop{\mathcal{E}\!\mathit{xt}}\nolimits}
\def\Sch{\mathit{Sch}}
\def\Mor{\mathop{\mathrm{Mor}}\nolimits}
\def\Ob{\mathop{\mathrm{Ob}}\nolimits}
\def\Sh{\mathop{\mathit{Sh}}\nolimits}
\def\NL{\mathop{N\!L}\nolimits}
\def\CH{\mathop{\mathrm{CH}}\nolimits}
\def\proetale{{pro\text{-}\acute{e}tale}}
\def\etale{{\acute{e}tale}}
\def\QCoh{\mathit{QCoh}}
\def\Ker{\mathop{\mathrm{Ker}}}
\def\Im{\mathop{\mathrm{Im}}}
\def\Coker{\mathop{\mathrm{Coker}}}
\def\Coim{\mathop{\mathrm{Coim}}}

% Boxtimes
%
\DeclareMathSymbol{\boxtimes}{\mathbin}{AMSa}{"02}

%
% Macros for moduli stacks/spaces
%
\def\QCohstack{\mathcal{QC}\!\mathit{oh}}
\def\Cohstack{\mathcal{C}\!\mathit{oh}}
\def\Spacesstack{\mathcal{S}\!\mathit{paces}}
\def\Quotfunctor{\mathrm{Quot}}
\def\Hilbfunctor{\mathrm{Hilb}}
\def\Curvesstack{\mathcal{C}\!\mathit{urves}}
\def\Polarizedstack{\mathcal{P}\!\mathit{olarized}}
\def\Complexesstack{\mathcal{C}\!\mathit{omplexes}}
% \Pic is the operator that assigns to X its picard group, usage \Pic(X)
% \Picardstack_{X/B} denotes the Picard stack of X over B
% \Picardfunctor_{X/B} denotes the Picard functor of X over B
\def\Pic{\mathop{\mathrm{Pic}}\nolimits}
\def\Picardstack{\mathcal{P}\!\mathit{ic}}
\def\Picardfunctor{\mathrm{Pic}}
\def\Deformationcategory{\mathcal{D}\!\mathit{ef}}
\usepackage{cleveref}
\crefname{originallemma}{original lemma}{original lemmas}
\Crefname{originallemma}{Original lemma}{Original lemmas}
\makeatletter
\newlabel{lemma-closed-immersion-ideals}{{01QP}{source}{}{sourcepointer-1}{}}
\newlabel{lemma-closed-immersion-ideals@cref}{{[originallemma][0][]01QP}{[1][0][]source}{}{}{}}
\makeatother
\crefname{originallemma}{original lemma}{original lemmas}
\Crefname{originallemma}{Original lemma}{Original lemmas}
\makeatletter
\newlabel{lemma-i-star-equivalence}{{01QY}{source}{}{sourcepointer-2}{}}
\newlabel{lemma-i-star-equivalence@cref}{{[originallemma][0][]01QY}{[1][0][]source}{}{}{}}
\makeatother
\crefname{originalsection}{original section}{original sections}
\Crefname{originalsection}{Original section}{Original sections}
\makeatletter
\newlabel{schemes-section-glueing-schemes}{{01JA}{source}{}{sourcepointer-3}{}}
\newlabel{schemes-section-glueing-schemes@cref}{{[originalsection][0][]01JA}{[1][0][]source}{}{}{}}
\makeatother
\title{The scheme theoretic union of two closed subschemes is the pushout of one along their scheme theoretic intersection, and its structure sheaf sits in a short exact sequence}
\author{The Stacks Project Authors}
\date{stacks-project:morphisms.tex:local-snapshot-20261012; GNU FDL 1.2 or later}
\begin{document}
\maketitle
\noindent\textit{Editorial note.} All delivered mathematical text is the author's own, copied from the pinned Stacks Project chapter file morphisms.tex; no statement or proof step was rewritten, reordered or removed.\par\medskip
\noindent\textit{Editorial note.} The canonical Stacks tag of the counted statement is 0C4J (morphisms-lemma-scheme-theoretic-union). Following the multi-claim convention for this source, source\_url is the source-level URL https://stacks.math.columbia.edu/ and the tag is carried in claim\_key and here rather than in a per-claim tag URL.\par\medskip
\noindent\textit{Editorial note.} The author's slip at proof line 610, where the restriction of g is written g|\_\{Y cap V\} while the surrounding argument uses Y cap U, is reproduced verbatim; no mathematical text was corrected.\par\medskip
\noindent\textit{Editorial note.} The preamble is the chapter's shared Stacks Project preamble, retained as sources/stacks-morphisms/retained-preamble.tex: a line-for-line excerpt of sources/stacks-chow/preamble.tex lines 14-18 and 154-217. No macro was invented or redefined.\par\medskip
\noindent\textit{Editorial note.} The definition of the scheme theoretic intersection and union is included as a context span because the statement names X union Y, X cap Y without repeating the definition.\par\medskip
\noindent\textit{Allowed standard prerequisites (editorial).}\begin{itemize}
\item Closed subschemes of a scheme correspond to quasi-coherent sheaves of ideals, and inclusion of subschemes is inclusion of ideals. Same chapter, Lemma labelled lemma-closed-immersion-ideals (canonical Stacks tag 01QP); invoked at proof line 577. Supplies that X -> X union Y and Y -> X union Y are closed immersions, and that the union is cut out by the intersection of the two ideal sheaves.
\item For a closed immersion i : Z -> X with ideal sheaf I, the functor i\_* identifies quasi-coherent O\_Z-modules with quasi-coherent O\_X-modules annihilated by I. Same chapter, Lemma labelled lemma-i-star-equivalence (canonical Stacks tag 01QY); invoked at proof line 578. Lets the author regard all four structure sheaves of the displayed sequence as quasi-coherent O\_S-modules, so that exactness may be checked on stalks or affine locally.
\item Morphisms of schemes may be constructed by glueing local data. Schemes, Section labelled schemes-section-glueing-schemes (canonical Stacks tag 01JA); invoked at proof line 601. Justifies building the morphism h : X union Y -> T from its restrictions to the three regions of the proof.
\end{itemize}\medskip
\ifdefined\hypertarget\hypertarget{sourcepointer-1}{}\fi
\noindent\textit{Source reference.} Original lemma 01QP: sources/stacks-morphisms/pinned-original-source.tex, source lines 110--110. Same-chapter pointer invoked at proof line 577 for the ideal-sheaf description of closed subschemes; the ideal sheaves I and J themselves are named in the counted definition and the conclusion that the union is cut out by I cap J is used inside the counted proof.\par\medskip
\ifdefined\hypertarget\hypertarget{sourcepointer-2}{}\fi
\noindent\textit{Source reference.} Original lemma 01QY: sources/stacks-morphisms/pinned-original-source.tex, source lines 386--386. Same-chapter pointer invoked at proof line 578 as a way of viewing the four structure sheaves as quasi-coherent O\_S-modules; all exactness and glueing arguments that use this point of view are carried out inside the counted proof.\par\medskip
\ifdefined\hypertarget\hypertarget{sourcepointer-3}{}\fi
\noindent\textit{Source reference.} Original section 01JA: sources/stacks-chow/tags.txt, source lines 2007--2007. Cross-chapter pointer invoked at proof line 601 for the licence to construct a morphism from local data; the local construction of h and the case analysis are carried out in full inside the counted span.\par\medskip
\begin{lemma}
\label{lemma-scheme-theoretic-union}
Let $S$ be a scheme. Let $X, Y \subset S$ be closed subschemes.
Let $X \cup Y$ be the scheme theoretic union of $X$ and $Y$.
Let $X \cap Y$ be the scheme theoretic intersection of $X$ and $Y$.
Then $X \to X \cup Y$ and $Y \to X \cup Y$ are closed immersions, there is a
short exact sequence
$$
0 \to \mathcal{O}_{X \cup Y} \to \mathcal{O}_X \times \mathcal{O}_Y
\to \mathcal{O}_{X \cap Y} \to 0
$$
of $\mathcal{O}_S$-modules, and the diagram
$$
\xymatrix{
X \cap Y \ar[r] \ar[d] & X \ar[d] \\
Y \ar[r] & X \cup Y
}
$$
is cocartesian in the category of schemes, i.e.,
$X \cup Y = X \amalg_{X \cap Y} Y$.
\end{lemma}

\begin{proof}
The morphisms $X \to X \cup Y$ and $Y \to X \cup Y$ are closed immersions
by Lemma \ref{lemma-closed-immersion-ideals}. In the short exact sequence
we use the equivalence of Lemma \ref{lemma-i-star-equivalence} to think of
quasi-coherent modules on closed subschemes of $S$ as quasi-coherent modules
on $S$. For the first map in the sequence we use the canonical maps
$\mathcal{O}_{X \cup Y} \to \mathcal{O}_X$ and
$\mathcal{O}_{X \cup Y} \to \mathcal{O}_Y$
and for the second map we use the canonical map
$\mathcal{O}_X \to \mathcal{O}_{X \cap Y}$ and
the negative of the canonical map
$\mathcal{O}_Y \to \mathcal{O}_{X \cap Y}$. Then to check
exactness we may work affine locally.
Let $U = \Spec(A)$ be an affine open of $S$ and let $X \cap U$ and $Y \cap U$
correspond to the ideals $I \subset A$ and $J \subset A$. Then
$(X \cup Y) \cap U$ corresponds to $I \cap J \subset A$
and $X \cap Y \cap U$ corresponds to $I + J \subset A$.
Thus exactness follows from the exactness of
$$
0 \to A/I \cap J \to A/I \times A/J \to A/(I + J) \to 0
$$
To show the diagram is cocartesian, suppose we are given a scheme $T$
and morphisms of schemes $f : X \to T$, $g : Y \to T$ agreeing
as morphisms $X \cap Y \to T$. Goal: Show there exists a unique
morphism $h : X \cup Y \to T$ agreeing with $f$ and $g$.
To construct $h$ we may work affine locally on $X \cup Y$, see
Schemes, Section \ref{schemes-section-glueing-schemes}.
If $s \in X$, $s \not \in Y$, then $X \to X \cup Y$ is
an isomorphism in a neighbourhood of $s$ and it is clear
how to construct $h$. Similarly for $s \in Y$, $s \not \in X$.
For $s \in X \cap Y$ we can pick an affine open
$V = \Spec(B) \subset T$ containing $f(s) = g(s)$.
Then we can choose an affine open $U = \Spec(A) \subset S$
containing $s$ such that $f(X \cap U)$ and $g(Y \cap U)$
are contained in $V$. The morphisms $f|_{X \cap U}$
and $g|_{Y \cap V}$ into $V$ correspond to ring maps
$$
B \to A/I
\quad\text{and}\quad
B \to A/J
$$
which agree as maps into $A/(I + J)$. By the short exact sequence
displayed above there is a unique lift of these ring homomorphism
to a ring map $B \to A/I \cap J$ as desired.
\end{proof}

\begin{definition}
\label{definition-scheme-theoretic-intersection-union}
Let $X$ be a scheme. Let $Z, Y \subset X$ be closed subschemes
corresponding to quasi-coherent ideal sheaves
$\mathcal{I}, \mathcal{J} \subset \mathcal{O}_X$.
The {\it scheme theoretic intersection} of $Z$ and $Y$
is the closed subscheme of $X$ cut out by $\mathcal{I} + \mathcal{J}$.
The {\it scheme theoretic union} of $Z$ and $Y$
is the closed subscheme of $X$ cut out by
$\mathcal{I} \cap \mathcal{J}$.
\end{definition}

\end{document}
